\chapter{Radiation Dosimetry and Contrast Safety Reference}
\label{app:radiation_contrast}

This appendix provides a quick reference for radiation dose metrics, contrast media screening, and MRI safety principles. Institutional policies and current ACR/IAEA guidance govern clinical protocols.

\section{Radiation Dose Metrics and Typical Adult Ranges}

\begin{longtable}{@{}p{0.34\textwidth}p{0.22\textwidth}p{0.36\textwidth}@{}}
\caption{Radiation dose quantities and representative adult effective doses.}
\label{tab:app_dose_reference} \\
\toprule
\textbf{Examination / Metric} & \textbf{Typical Range} & \textbf{Notes} \\
\midrule
\endfirsthead
\toprule
\textbf{Examination / Metric} & \textbf{Typical Range} & \textbf{Notes} \\
\midrule
\endhead
Annual natural background & $\approx 3\text{ mSv/year}$ & Varies by geography and altitude \\
Chest radiograph (PA/lateral) & 0.02--0.1 mSv & Very low dose \\
CT head (non-contrast) & 1.5--2.5 mSv & Output reported as $\text{CTDI}_{\text{vol}}$ and DLP \\
CT chest & 5--8 mSv & Low-dose lung screening $\approx 1.0\text{--}1.5\text{ mSv}$ \\
CT abdomen and pelvis & 7--12 mSv & Multiphase protocols increase dose proportionally \\
Whole-body FDG PET/CT & 14--25 mSv & Combines radiotracer and CT components \\
\bottomrule
\end{longtable}

\section{Contrast Media Safety Summary}

\begin{longtable}{@{}p{0.22\textwidth}p{0.34\textwidth}p{0.36\textwidth}@{}}
\caption{Key safety considerations for iodinated and gadolinium-based contrast agents.}
\label{tab:app_contrast_safety} \\
\toprule
\textbf{Topic} & \textbf{Iodinated Contrast (CT/Angio)} & \textbf{Gadolinium-Based Contrast (MRI)} \\
\midrule
\endfirsthead
\toprule
\textbf{Topic} & \textbf{Iodinated Contrast (CT/Angio)} & \textbf{Gadolinium-Based Contrast (MRI)} \\
\midrule
\endhead
Prior reaction & Screen for prior moderate/severe reaction; consider alternative modality or premedication per policy & Switch agent class or use non-contrast MRI if prior reaction \\
Seafood/shellfish allergy & Not a specific predictor of iodinated contrast reaction & Not relevant \\
Kidney function & Assess eGFR per local threshold; hydrate high-risk patients & Use Group II macrocyclic agents when clinically indicated; negligible NSF risk \\
Pregnancy & Use only when clinically justified; no routine neonatal thyroid screening beyond standard newborn screen & Avoid unless essential; weigh maternal/fetal benefit \\
Breastfeeding & Safe to continue breastfeeding without interruption & Safe to continue breastfeeding with standard macrocyclic agents \\
\bottomrule
\end{longtable}

\section{MRI Safety Zones}
\begin{itemize}
    \item \textbf{Zone I:} General public access; no magnetic hazard.
    \item \textbf{Zone II:} Patient reception and MRI safety screening interface.
    \item \textbf{Zone III:} Controlled access area near the magnet room; ferromagnetic objects are strictly restricted.
    \item \textbf{Zone IV:} The magnet room itself (static $B_0$ field always active); only MR-safe or verified MR-conditional equipment permitted.
\end{itemize}
